\subsubsection{Election result files (\texttt{elections/CC/\ldots.csv})}
\label{sec:variables_election_files}

\begin{longtable}{p{3cm}p{12cm}}
\toprule
\textbf{Variable} & \textbf{Description} \\
\midrule
\endhead
\texttt{gisco\_id}      & Geographic unit identifier.  Links to
                          \texttt{geo/laus.csv}, \texttt{geo/nuts.csv} or
                          \texttt{geo/special.csv}. \\
\texttt{name}           & Canonical geographic unit name valid in the
                          election year (resolved from \texttt{name\_or}
                          and historical name variants in the geo files). \\
\texttt{registered}     & Number of registered voters (integer). \\
\texttt{turnout}        & Total votes cast (valid and invalid combined) (integer). \\
\texttt{invalid}        & Number of invalid or spoiled ballots (integer). \\
\texttt{blank}          & Number of blank ballots (integer). \\
\textit{party columns}  & One column per distinct \texttt{party\_id} used in
                          the country file.
                          Values are absolute vote counts (integers). \\
\bottomrule
\end{longtable}

Details about specific vote counts and missing data at the municipal level
can be found in the country-specific notes.

\paragraph{Party column naming.}
Party columns use \texttt{party\_id} values from
\texttt{parties/parties.csv}.  The \texttt{party\_id} key is unique within
each country and resolves source-file aliases to a single party row.

\paragraph{File naming convention.}
Election files follow the convention
$$\texttt{\{year\}\_\{month\}[\_\{round\}][\_\{vote\_type\}].csv}$$
for national elections and
$$\texttt{EU\_\{year\}\_\{month\}.csv}$$
for European Parliament elections, where \textit{month} is the zero-padded
two-digit month (e.g.\ \texttt{09} for September).
The optional \textit{round} suffix distinguishes voting rounds in two-round
systems (e.g.\ \texttt{2024\_06\_1.csv} and \texttt{2024\_07\_2.csv} for
France~2024).
The optional \textit{vote\_type} suffix distinguishes ballot types where a
single election produces separate tallies (e.g.\
\texttt{Erststimme} or \texttt{Zweitstimme} for Germany,
\texttt{party} or \texttt{single\_member} for Hungary).

\subsubsection{Election index (\texttt{data\_files\_reference.csv})}
\label{sec:variables_reference}

A flat index of all election result files included in the distribution.

\begin{longtable}{p{3.5cm}p{11.5cm}}
\toprule
\textbf{Variable} & \textbf{Description} \\
\midrule
\endhead
\texttt{country\_code}    & ISO~3166 country code (e.g.\ \texttt{AT},
                            \texttt{DE}). \\
\texttt{parliament\_name} & Name of the parliament or assembly in the
                            national language. \\
\texttt{election\_date}   & Date of the election (\texttt{YYYY-MM-DD}).
                            For two-round elections this is the date of the
                            specific round; for multi-day elections it is
                            the last polling day. \\
\texttt{round}            & Round number (integer, empty for single-round
                            elections). \\
\texttt{vote\_type}       & Ballot type label for elections with multiple
                            simultaneous tallies (e.g.\ \texttt{Erststimme},
                            \texttt{party}); empty otherwise. \\
\texttt{file\_path}       & Path to the result file relative to the
                            distribution root (e.g.\
                            \texttt{elections/AT/2002\_11.csv}). \\
\bottomrule
\end{longtable}
